\documentclass[11pt,letterpaper]{article}

% ── Geometry ──────────────────────────────────────────────────────────────────
\usepackage[margin=1in]{geometry}

% ── Pagination ────────────────────────────────────────────────────────────────
\usepackage[all]{nowidow}    % Require ≥2 lines on each side of a page break
\brokenpenalty=10000         % No hyphenated words across page breaks
\raggedbottom                % Allow uneven page bottoms rather than stretching

% ── Fonts ─────────────────────────────────────────────────────────────────────
\usepackage[T1]{fontenc}
\usepackage{newpxtext}     % Palatino-based serif (TeX Gyre Pagella)
\usepackage{newpxmath}     % Matching math font

% ── Colors ────────────────────────────────────────────────────────────────────
\usepackage[dvipsnames]{xcolor}
\definecolor{SectionBlue}{HTML}{1B3A5C}
\definecolor{AccentGray}{HTML}{4A4A4A}
\definecolor{QuoteBorder}{HTML}{8B9DAF}
\definecolor{RiskRed}{HTML}{C0392B}
\definecolor{RiskAmber}{HTML}{E67E22}
\definecolor{RiskGreen}{HTML}{27AE60}
\definecolor{BoxBg}{HTML}{F7F8FA}

% ── Tables ────────────────────────────────────────────────────────────────────
\usepackage{booktabs}
\usepackage{tabularx}

% ── Boxes ─────────────────────────────────────────────────────────────────────
\usepackage[framemethod=tikz]{mdframed}

% ── Lists ─────────────────────────────────────────────────────────────────────
\usepackage{enumitem}
\usepackage{needspace}

% ── Hyperlinks ────────────────────────────────────────────────────────────────
\usepackage{hyperref}
\hypersetup{
  colorlinks=true,
  linkcolor=SectionBlue,
  urlcolor=SectionBlue,
  citecolor=SectionBlue,
  pdfauthor={Punt Labs},
  pdftitle={Lux v2 PR/FAQ},
  pdfsubject={Working Backwards — Agent IDE},
}

% ── Bibliography ─────────────────────────────────────────────────────────────
\usepackage[style=numeric,backend=biber,sorting=none]{biblatex}
\addbibresource{lux-prfaq-v2.bib}

% ── Section styling ───────────────────────────────────────────────────────────
\usepackage{titlesec}
\titleformat{\section}
  {\needspace{6\baselineskip}\Large\bfseries\color{SectionBlue}}
  {}
  {0pt}
  {}

\titleformat{\subsection}
  {\needspace{5\baselineskip}\large\bfseries\color{AccentGray}}
  {}
  {0pt}
  {}

\titleformat{\subsubsection}
  {\normalsize\bfseries\color{AccentGray}}
  {}
  {0pt}
  {}
\titlespacing*{\subsubsection}{0pt}{1em}{0.5em}[0pt]

% ── Document stage ───────────────────────────────────────────────────────────
\newcommand{\prfaqstage}[1]{\def\theprfaqstage{#1}}
\prfaqstage{hypothesis}

% ── Document version ─────────────────────────────────────────────────────────
\newcommand{\prfaqversion}[2]{\def\theprfaqmajor{#1}\def\theprfaqminor{#2}}
\prfaqversion{1}{0}

% ── Header/footer ────────────────────────────────────────────────────────────
\usepackage{fancyhdr}
\pagestyle{fancy}
\fancyhf{}
\fancyhead[L]{\footnotesize\textcolor{AccentGray}{Stage: \theprfaqstage{} \enspace\textbar\enspace v\theprfaqmajor.\theprfaqminor}}
\fancyhead[R]{\footnotesize\textcolor{AccentGray}{\thepage}}
\renewcommand{\headrulewidth}{0pt}

% ══════════════════════════════════════════════════════════════════════════════
% Custom environments
% ══════════════════════════════════════════════════════════════════════════════

\newcommand{\prsection}[1]{%
  \vspace{1em}%
  {\large\bfseries\color{SectionBlue}#1}%
  \vspace{0.3em}\par%
}

\newmdenv[
  topline=false,
  bottomline=false,
  rightline=false,
  linewidth=3pt,
  linecolor=QuoteBorder,
  backgroundcolor=BoxBg,
  innerleftmargin=12pt,
  innerrightmargin=12pt,
  innertopmargin=8pt,
  innerbottommargin=8pt,
  skipabove=\baselineskip,
  skipbelow=\baselineskip,
]{customerquote}

\newmdenv[
  topline=false,
  bottomline=false,
  rightline=false,
  linewidth=3pt,
  linecolor=SectionBlue,
  backgroundcolor=BoxBg,
  innerleftmargin=12pt,
  innerrightmargin=12pt,
  innertopmargin=8pt,
  innerbottommargin=8pt,
  skipabove=\baselineskip,
  skipbelow=\baselineskip,
]{spokespersonquote}

\newcounter{faqnum}
\newenvironment{faqpair}[1]{%
  \needspace{4\baselineskip}%
  \vspace{0.5em}%
  \refstepcounter{faqnum}%
  \noindent\textbf{\color{SectionBlue}Q\thefaqnum: #1}\par\nopagebreak[4]%
  \vspace{0.2em}%
  \begin{adjustwidth}{1em}{0pt}%
  \setlength{\parindent}{0pt}%
}{%
  \end{adjustwidth}%
  \vspace{0.5em}%
}
\newcommand{\faqref}[1]{\hyperref[#1]{FAQ~\ref*{#1}}}

\newcounter{featurenum}
\newcommand{\featureitem}[2]{%
  \refstepcounter{featurenum}%
  \item[\textbf{\color{SectionBlue}F\thefeaturenum.}]%
  \textbf{#1} --- #2%
}
\newcommand{\featureref}[1]{\hyperref[#1]{Feature~\ref*{#1}}}

\usepackage{changepage}

% ══════════════════════════════════════════════════════════════════════════════
% Document
% ══════════════════════════════════════════════════════════════════════════════

\begin{document}

% ── Headline block ───────────────────────────────────────────────────────────
\begin{center}
  {\LARGE\bfseries\color{SectionBlue}
    Self-Extending Display Compositor Gives AI Agents\\
    a Shared Visual Surface That Grows On Demand} \\[0.8em]
  {\large\color{AccentGray}
    Lux~v2 lets agents author new visual capabilities in minutes ---
    no vendor release required, no web stack, pure Python} \\[0.3em]
  {\large\color{AccentGray}
    Open-source compositor with community extension registry,
    three rendering modes, and an in-display agent SDK}
\end{center}

\vspace{0.5em}

% ══════════════════════════════════════════════════════════════════════════════
% PRESS RELEASE
% ══════════════════════════════════════════════════════════════════════════════

\section*{Press Release}

AUSTIN, TX --- October 2026 --- Punt Labs today released Lux~v2, a
Python-native display compositor that gives AI coding agents a persistent,
shared visual surface.  Unlike terminal scrollback or vendor dashboards,
Lux~v2 extends itself: when an agent needs a visualization that does not
exist, it authors a new extension, installs it, and renders --- in under
five minutes, without a product release.  Extensions are shareable through
a community registry with ratings and verification badges, so most users
never need to author code themselves
(see \faqref{faq:community}).

\prsection{Problem}

Developers running AI coding agents --- Claude Code\cite{anthropic2025claudecode},
Cursor, Codex --- coordinate multiple agents across repos, sessions, and
machines.  The only visibility into what those agents are doing is terminal
scrollback.  As Steve Yegge has noted, coordinating multiple agents is ``hard
and exhausting''\cite{yegge2025agents}.  When an agent produces structured
output --- a table, a chart, an architecture diagram --- the developer sees
raw text or takes an ad~hoc screenshot.

Every vendor is building a web-based solution: dashboards, chat panels,
embedded notebooks.  These force the developer to wait for the vendor to
ship the specific view they need (see \faqref{faq:competition}).
The developer's terminal environment has no structured visual output.

\prsection{Solution}

Lux~v2 is a persistent local service --- like Ollama\cite{ollama2025} for
inference, but for display.  Agents connect via an HTTP service API and send
structured content.  The display renders it in a native window at 60\,fps
using Dear~ImGui\cite{imgui2025}.  Multiple agents share one display, each
in its own frame, with clear attribution.

The core handles transport, protocol dispatch, extension registry, frame
management, render loop, and safety.  Everything the user sees ---
every element kind, every applet, every theme --- is an extension.  Lux ships
with a starter pack of bundled extensions (text, tables, plots, diagrams,
interactive controls), but the vocabulary is not fixed.

When no installed extension covers a visualization need, an agent can author
one: generate the manifest and renderer code, run tests, and install ---
all through the service API.  The user approves via a consent dialog.
The extension persists in the local registry and is available for every
future session.  Publishing to the community registry lets anyone install it
with one command (see \faqref{faq:authoring}).

Lux is not a passive display.  The bidirectional write path --- user and
agent both driving the display, the display driving agents back --- was the
key lesson from v1 (see \faqref{faq:v1-lessons}).  An in-display agent
built on the Anthropic Python SDK can receive intent from external agents,
author extensions, introspect frames, and take screenshots for iterative GUI
refinement --- a technique proven during the Postern/Pharo
work\cite{postern2026,claudeagentsdk2026} (see \faqref{faq:pharo}).

\prsection{Customer Quote}

\begin{customerquote}
\textit{``I run six Claude Code sessions across three repos.  Before Lux, I
had six terminal windows and no idea what any of them were doing unless I
scrolled back.  Now I glance at one window and see beads boards, test
dashboards, and architecture diagrams --- one frame per agent, all live.
Last week an agent needed a dependency graph and there was no built-in
widget for it.  It wrote one.  Took three minutes.  I've been using that
widget every day since.''}

\raggedleft --- \textbf{Maya Torres}, Staff Engineer at a Series~B DevTools Startup
\end{customerquote}

\prsection{Getting Started}

Getting started takes one command:

\begin{enumerate}[nosep]
  \item Install:

    {\small\texttt{curl -fsSL https://raw.githubusercontent.com/\allowbreak punt-labs/lux/main/install.sh | sh}}
  \item Restart Claude Code.  The Lux display opens automatically when any
    agent sends visual output.
  \item Browse community extensions:
    {\small\texttt{lux ext search <keyword>}} and install with
    {\small\texttt{lux ext install <name>}}.
\end{enumerate}

No API key is required for the display or community extensions.  The
optional in-display agent (for authoring new extensions) requires an
Anthropic API key.

\prsection{Spokesperson Quote}

\begin{spokespersonquote}
``Lux~v1 proved the display has value --- the beads browser alone changed
how we work.  Pharo proved a self-extending system is valuable --- an agent
that writes its own UI code, tests it, and persists it.  Lux~v2 combines
both: the pragmatic Python-native display from v1, plus the self-extension
model from Pharo, grounded in Python's ecosystem.  The result is that
`Lux can't do X' stops being a limitation and becomes a five-minute
problem.''

\raggedleft --- \textbf{Jim Freeman}, Punt Labs
\end{spokespersonquote}

\prsection{Call to Action}

Lux~v2 is free and open-source under the MIT license.  Visit
{\small\texttt{https://github.com/punt-labs/lux}} to install, browse the
extension registry, or contribute.  The community registry is open for
submissions.

\newpage

% ══════════════════════════════════════════════════════════════════════════════
% FREQUENTLY ASKED QUESTIONS
% ══════════════════════════════════════════════════════════════════════════════

\section*{Frequently Asked Questions}

% ── External FAQs (Customer-facing) ──────────────────────────────────────────

\subsection*{External FAQs}

\begin{faqpair}{What is Lux and who is it for?}\label{faq:what}
  Lux is a Python-native display compositor --- a persistent local visual
  surface that AI coding agents write to and read from.  It is for developers who
  use CLI-based coding agents (Claude Code, Cursor, Codex, or similar) and
  need structured visual output: tables, charts, diagrams, interactive
  controls, and custom visualizations.  Lux runs on macOS and Linux.
\end{faqpair}

\begin{faqpair}{How is this different from terminal scrollback, Claude Desktop, or vendor dashboards?}\label{faq:competition}
  Terminal scrollback is unstructured text.  Claude Desktop and vendor
  dashboards are structured but closed: you get the views the vendor ships,
  on the vendor's release schedule.  Lux is open and extensible: when you
  need a visualization that does not exist, an agent authors it in minutes
  and it persists for every future session.

  Lux is also Python-native, not web.  It integrates with Python's
  visualization ecosystem (matplotlib, plotly, PIL, pyte, Qt, tkinter)
  directly, without transpiling to JavaScript.
\end{faqpair}

\begin{faqpair}{How do I get started?}
  One curl command installs Lux and registers the Claude Code plugin.
  Restart Claude Code.  The display opens automatically when any agent
  sends visual output.  Most users see their first visualization within
  60 seconds of install.
\end{faqpair}

\begin{faqpair}{How much does it cost?}
  Free and open-source (MIT).  Community extensions are free.  The optional
  in-display agent for authoring new extensions requires an Anthropic API
  key (usage-based, same pricing as any Claude API call).
\end{faqpair}

\begin{faqpair}{What happens to my data?}
  All data stays local.  Lux communicates over a Unix domain socket
  (localhost only by default) or TCP with TLS for remote agents.  No
  telemetry, no data leaves your machine.  Extensions run in-process;
  the consent model gates what code executes.
\end{faqpair}

\begin{faqpair}{Can I use Lux without enabling LLM-powered extension authoring?}
\label{faq:community}
  Yes.  The community registry provides pre-built extensions authored,
  tested, and rated by other users.  Install with
  {\small\texttt{lux ext install <name>}}.  No API key required.  No LLM
  runs in your Lux instance.  The extension store shows verification
  badges --- lint passing, CI green, test coverage, community ratings ---
  so you can assess quality before installing.

  A small number of extension authors produce extensions consumed by a
  larger population of users who never need to write code.
\end{faqpair}

\begin{faqpair}{What Python visualization libraries can Lux host?}
  Lux supports three rendering modes.  \textbf{ImGui mode}: extensions
  render natively in the Lux window using Dear~ImGui --- full interactivity,
  single-window UX.  \textbf{Texture mode}: extensions produce pixel
  buffers (matplotlib, PIL, plotly snapshots) displayed as images in the Lux
  window --- limited interactivity but zero integration effort.
  \textbf{Window mode}: extensions own their own OS window (Qt, tkinter,
  Jupyter kernels) --- full interactivity, coordinated lifecycle, separate
  window.

  Any Python library that can render to a buffer or open a window works.
\end{faqpair}

% ── Internal FAQs (Business-facing) ──────────────────────────────────────────

\newpage
\subsection*{Internal FAQs}

\subsubsection*{Value \& Market}

\begin{faqpair}{What is the total addressable market?}\label{faq:tam}
  The primary market is developers using CLI-based AI coding agents.
  Claude Code, Cursor (CLI mode), and Codex are the leading tools.
  Anthropic reports millions of Claude Code users.  The secondary market
  is any developer running Python-based AI pipelines who needs visual
  output.

  Bottoms-up: 1\% adoption of active Claude Code users (free, zero-config
  install) is tens of thousands of active installations.  Registry value
  scales with adoption --- more users means more extensions, which means
  more users.  [CITATION NEEDED: Claude Code user count]
\end{faqpair}

\begin{faqpair}{What evidence do we have that customers want this?}\label{faq:evidence}
  Lux~v1 shipped in March 2026 and has been in daily use across Punt Labs
  projects since.  The beads issue browser --- a filterable table that
  auto-refreshes on every \texttt{bd} command --- runs in every active
  agent session.  v1 proved three things: (1)~the display has value,
  (2)~the hook architecture (PostToolUse auto-refresh) makes it feel live,
  and (3)~agents and humans share one display surface effectively.

  The self-extension hypothesis is validated by the Postern/Pharo
  work\cite{postern2026}: 199~classes, 2,358~methods, and 765~tests
  written entirely by agents driving a live image via HTTP --- authored UI
  code, screenshotted, iterated, committed, all without human GUI
  interaction.

  What v1 did \emph{not} prove: whether community sharing works, whether
  the three rendering modes compose well, and whether the in-display agent
  is useful outside Punt Labs.  These are the hypotheses v2 must test.
\end{faqpair}

\begin{faqpair}{Who are the competitors and why will we win?}\label{faq:competitors}
  Every AI coding tool vendor is building visual output for agents:

  \begin{itemize}[nosep]
    \item \textbf{Claude Desktop} --- Anthropic's consumer product.
      Web-based, integrated with the chat interface.  Closed; users
      cannot add custom visualizations.
    \item \textbf{Cursor} --- IDE with AI features.  Visual output is
      in-editor.  Extensible via VS Code extensions, but these are
      IDE-scoped and written by hand.
    \item \textbf{Codex} --- OpenAI's coding agent.  Terminal-only as of
      2026.  No visual output surface.
    \item \textbf{Jupyter notebooks} --- Rich visual output via
      kernel/display protocol.  Not agent-first; requires notebook UX.
  \end{itemize}

  \textbf{Primary threat: Anthropic-native visual output.}  Anthropic has
  every incentive to build visual panels into Claude Code.  A basic
  built-in display narrows Lux's install-base advantage.  An extensible
  one collapses Lux's differentiation to ``Python-native vs.\ web-native.''
  Survival strategies: (a)~Lux's extension ecosystem reaches critical mass
  before Anthropic ships, becoming the de-facto standard; (b)~Lux becomes
  a complementary extension provider for whatever Anthropic builds;
  (c)~Python-native rendering and offline operation remain valuable for
  users who want local-only display.

  Lux's structural advantage: it grows its vocabulary via LLM-authored
  extensions.  Vendor products ship features on their schedule; Lux
  delivers a new capability in minutes because the agent writes it.
  A competitor copying the extension registry concept would have to abandon
  web-only rendering to integrate Python's visualization libraries ---
  a platform-level decision they are unlikely to reverse.
\end{faqpair}

\begin{faqpair}{What is your next step to validate the vision?}\label{faq:next-step}
  Two parallel tracks.

  \textbf{External demand:} Interview 5--10 developers running 3+
  concurrent agent sessions daily (staff-level, startups or mid-size
  companies).  Key questions: (a)~how do you track what your agents are
  doing? (b)~when an agent produces structured output, how do you consume
  it? (c)~would you install a display tool if it took 60 seconds?
  Go/no-go: if fewer than half describe agent visibility as a top-3 pain
  point, revisit the value proposition before committing to phases 4--6.

  \textbf{Technical:} Build the extension loader and registry as the first
  v2 milestone.  Migrate existing element kinds to registered extensions.
  Ship 3--5 reference extensions (matplotlib plot, terminal emulator,
  sankey diagram).  Measure: can an agent author and install a new
  extension in under five minutes?  Run the loop 10 times and report
  median and range.
\end{faqpair}

\subsubsection*{Technical}

\begin{faqpair}{What are the major technical risks?}\label{faq:tech-risks}
  \begin{enumerate}[nosep]
    \item \textbf{Extension safety.}  Python has no in-process sandbox;
      extensions run with full process access.  Mitigation: exception
      isolation (crash $\to$ error placeholder, core keeps rendering),
      consent dialog on install, community verification badges (lint, CI,
      tests).  The trust model is identical to \texttt{pip install} ---
      the user is the security boundary.
    \item \textbf{Python dynamism limits.}  Python cannot do Smalltalk's
      \texttt{become} or image persistence.  Extensions persist as files;
      the process is disposable.  Adequate for hot-loading new element
      kinds, but the display cannot save arbitrary runtime state across
      restarts.
    \item \textbf{Three rendering modes composing well.}  ImGui, texture,
      and window modes coexist in one session.  The risk is UX
      fragmentation --- separate OS windows feel disconnected.  Mitigation:
      prefer ImGui mode; texture mode for most Python viz libraries;
      window mode as escape hatch only.
    \item \textbf{Community registry bootstrapping.}  A registry with no
      extensions has no users.  Mitigation: ship 20+ starter-pack
      extensions plus 5+ reference extensions from day one.
  \end{enumerate}
\end{faqpair}

\begin{faqpair}{What is the estimated development timeline?}\label{faq:timeline}
  Phased delivery, ordered by dependency.  Scope compresses by cutting
  later phases.

  \begin{enumerate}[nosep]
    \item \textbf{Extension system.}  Registry, loader, manifest format,
      safety wrapper.  Migrate existing element kinds to extensions.
    \item \textbf{Service API.}  HTTP endpoints (ollama-shaped).  MCP
      adapter becomes a thin translation layer.  Auto-documenting
      \texttt{/help} endpoint.
    \item \textbf{Reference extensions.}  Matplotlib, terminal (pyte),
      plotly snapshot, custom ImGui widgets.  Proves the extension
      contract with real libraries.
    \item \textbf{In-display agent.}  Anthropic SDK integration.
      LuxExchange tool-use loop.  Leaf and composite commands.
      Permission policy.  Screenshot-driven iteration.
    \item \textbf{Community registry.}  GitHub-hosted.  CLI:
      \texttt{lux ext install/publish/search}.  Verification badges.
      Ratings.
    \item \textbf{Remote transport.}  TCP with TLS.  Multi-machine
      agent sessions.
  \end{enumerate}

  Cut line: phases 5--6 are valuable but not launch-blocking.  Core
  value is in phases 1--3.
\end{faqpair}

\begin{faqpair}{What is the relationship to Pharo and Postern?}\label{faq:pharo}
  Complementary, not competing.  Pharo is the ambitious live-environment
  substrate: full image persistence, introspection, Morphic rendering.
  Postern\cite{postern2026} exposes Pharo over HTTP.  The Claude Agent
  SDK for Smalltalk\cite{claudeagentsdk2026} provides the in-image tool
  runner.

  Lux is the pragmatic Python-native path.  It shares the self-extension
  philosophy --- agent authors code, code persists, system grows --- but
  uses Python's ecosystem instead of Smalltalk's.

  The two connect: a Lux window-mode extension can host Pharo-rendered
  content via Postern, and extensions designed in Morphic can be exported
  as Lux extensions for Python consumers.
\end{faqpair}

\subsubsection*{Business}

\begin{faqpair}{What is the revenue model?}
  Lux is free and open-source (MIT).  Follow-on revenue opportunities:
  hosted managed instances, premium verified extensions, enterprise
  support.  None are launch requirements.

  Strategic value: Lux drives adoption of other Punt Labs tools via two
  mechanisms.  First, Lux surfaces those tools in its UI --- the beads
  browser, biff presence panels, vox controls, and quarry search all
  render as Lux extensions; users who see them discover the underlying
  tools.  Second, the install script bundles the full Punt Labs CLI suite.
  Validation target: 10\% of Lux users install one additional Punt Labs
  tool within 6~months.
\end{faqpair}

\begin{faqpair}{What are the key metrics for success?}\label{faq:metrics}
  \begin{enumerate}[nosep]
    \item \textbf{Daily active display sessions} --- target: 100 within
      3~months of v2 launch.
    \item \textbf{Community extensions published} --- target: 50 within
      6~months.
    \item \textbf{Extension install rate} --- target: average user
      installs 3+ community extensions.
    \item \textbf{Agent-authored extension success rate} --- target:
      80\% of authoring attempts produce a working extension on first
      try.
    \item \textbf{Time to new capability} --- target: median under
      5~minutes from ``I need X'' to ``X is rendering.''
  \end{enumerate}

  Decision threshold: if daily active sessions are below 30 after
  6~months, revisit whether the self-extending compositor framing
  resonates or whether Lux should remain a simpler display surface
  with a fixed vocabulary.
\end{faqpair}

\begin{faqpair}{Why now?  What has changed?}
  Three things converged.  (1)~AI coding agents are mainstream ---
  Claude Code, Cursor, and Codex have millions of users [CITATION NEEDED],
  all producing output invisible beyond terminal scrollback.  (2)~LLMs
  can write UI extensions on demand --- the authoring loop that was
  impossible two years ago is now routine (proven on Postern/Pharo).
  (3)~Lux~v1 proved the display surface has value and the
  protocol-as-API design works.
\end{faqpair}

\begin{faqpair}{What lessons did Lux~v1 teach?}\label{faq:v1-lessons}
  Five lessons from nine weeks of daily use:

  \begin{enumerate}[nosep]
    \item The display has value.  The beads browser changed daily workflow.
    \item Bidirectional matters.  Passive display is not enough; the
      display must drive agent behavior (via events, menu clicks,
      interaction messages).
    \item Fixed vocabularies are limiting.  27~element kinds is both too
      many (most unused) and too few (missing what users actually need).
    \item The hook architecture works.  PostToolUse auto-refresh makes
      the display feel live without explicit agent coordination.
    \item Multi-agent display is real.  Per-project frames, shared
      ownership, and orphan persistence solve the N-to-1 problem at the
      protocol level.
  \end{enumerate}
\end{faqpair}

\begin{faqpair}{How does the LLM extension authoring loop work?}\label{faq:authoring}
  An agent calls \texttt{/api/extensions/author} with a description.
  The in-display agent (Anthropic Python SDK) generates the extension:
  a manifest declaring the element kind, schema, dependencies, and
  rendering mode, plus a renderer function.  It runs tests.  A consent
  dialog shows the generated code.  The user approves.  The extension
  installs to \texttt{\textasciitilde/.lux/extensions/} and is available
  immediately.

  The in-display agent can screenshot the rendered result and iterate ---
  ``the table is misaligned, fix the column widths'' --- using the same
  technique proven on Postern/Pharo.  This visual feedback loop is what
  makes rapid authoring realistic.
\end{faqpair}

\begin{faqpair}{What resources does this require?}
  Lux~v2 is built by one engineer with AI agent assistance (Claude Code
  driving implementation, review, and testing --- the same model used for
  Lux~v1's 629 tests and the Smalltalk SDK's 811 tests).

  Phases 1--3 (extension system, service API, reference extensions) are
  the priority; other Punt Labs tools (biff, vox, quarry) slow to
  maintenance-only during that work.  Phases 4--6 are sequenced after
  phases 1--3 ship and external demand is validated.
\end{faqpair}

\begin{faqpair}{Pre-mortem: it is 12 months from now and Lux~v2 has
  failed.  What happened?}\label{faq:premortem}
  Five plausible failure scenarios:

  \begin{enumerate}[nosep]
    \item \textbf{Anthropic ships built-in visual output in Claude Code.}
      A first-party surface, even a basic one, eliminates Lux's
      install-friction advantage.  If it is also extensible, Lux's
      core differentiation collapses.
    \item \textbf{External developers do not care about agent visibility.}
      The Punt Labs team uses the beads browser daily, but this may be a
      niche workflow.  Fewer than half of interviewed developers describing
      visibility as a top-3 pain point means the market is smaller than
      projected.
    \item \textbf{Extension authoring success rate is below 50\%.}
      Agents failing more often than they succeed turns the ``five-minute''
      promise into ``thirty minutes of debugging,'' and users give up.
    \item \textbf{Community registry never bootstraps.}  Below 100 active
      users, the registry has too few extensions to attract users and too
      few users to attract authors.  The flywheel never spins.
    \item \textbf{Three rendering modes fragment the UX.}  Window-mode
      extensions feel disconnected; texture-mode interactivity is too
      limited.  Users perceive Lux as inconsistent rather than flexible.
  \end{enumerate}
\end{faqpair}

% ══════════════════════════════════════════════════════════════════════════════
% FOUR RISKS ASSESSMENT
% ══════════════════════════════════════════════════════════════════════════════

\newpage
\section*{Risk Assessment}

\renewcommand{\arraystretch}{1.6}
\begin{tabularx}{\textwidth}{@{} l l X @{}}
\toprule
\textbf{\color{SectionBlue}Risk} & \textbf{\color{SectionBlue}Rating} & \textbf{\color{SectionBlue}Assessment} \\
\midrule
Value      & \textcolor{RiskAmber}{Medium} &
  v1 proved the display has value for Punt Labs.  Unproven: external
  developer adoption, whether the extension model resonates beyond power
  users, and whether the agent IDE framing is the right positioning. \\
Usability  & \textcolor{RiskAmber}{Medium} &
  Install and basic display are proven (one curl command, auto-open).
  Extension consumption is plausible but untested with external users.
  Extension authoring --- the headline differentiator --- requires a
  multi-step consent flow (generate, review, approve, iterate via
  screenshot) not yet tested outside Punt Labs. \\
Feasibility & \textcolor{RiskAmber}{Medium} &
  Core rendering and IPC are proven.  Extension loader is a
  well-understood pattern.  New risks: three rendering modes composing
  coherently (only ImGui exists in v1), in-display agent SDK integration,
  community registry bootstrapping, and remote TCP/TLS.  No single risk
  is High, but the aggregate --- 12+ HTTP endpoints, a tool-use loop,
  five permission modes, three render contexts --- is substantial new
  work. \\
Viability  & \textcolor{RiskAmber}{Medium} &
  Open-source with no direct revenue.  Strategic value depends on
  ecosystem pull for other Punt Labs tools.  If Lux~v2 does not drive
  adoption of biff/vox/quarry/beadle, the investment is a public good
  without business return.  Acceptable if treated as infrastructure,
  not a profit center. \\
\bottomrule
\end{tabularx}
\renewcommand{\arraystretch}{1.0}

% ══════════════════════════════════════════════════════════════════════════════
% FEATURE APPENDIX
% ══════════════════════════════════════════════════════════════════════════════

\newpage
\section*{Feature Appendix}

Defines the scope boundary for Lux~v2.  Every feature is categorized
to prevent scope creep and force prioritization upfront.

\subsection*{Must Do}

Essential for launch.  Without these, Lux~v2 is not distinguishable
from v1.

\begin{enumerate}[nosep,leftmargin=2.5em]
  \featureitem{Extension registry and loader}{runtime dispatch table
    replacing hardcoded element kinds; scan, validate, register,
    crash-isolate}\label{feat:registry}
  \featureitem{Extension manifest format}{declared element kind, schema,
    dependencies, rendering mode (imgui/texture/window), metadata}\label{feat:manifest}
  \featureitem{Starter pack migration}{existing 27 element kinds
    refactored as bundled extensions, not core code}\label{feat:migration}
  \featureitem{HTTP service API}{ollama-shaped endpoints for show,
    update, events, extensions, capabilities; MCP becomes thin
    adapter}\label{feat:service-api}
  \featureitem{Auto-documenting \texttt{/help} endpoint}{runtime
    truth from live registry: installed kinds, services, conventions,
    permission mode}\label{feat:help}
  \featureitem{Extension CLI}{install, list, and remove extensions via
    \texttt{lux ext}; user-local persistence at
    \texttt{\char`\~/.lux/extensions/}}\label{feat:ext-cli}
  \featureitem{Safety wrapper}{exception isolation per extension
    render call; crash shows error placeholder, core keeps
    rendering}\label{feat:safety}
  \featureitem{Consent model for extension install}{show source,
    user approves; same trust model as render\_function}\label{feat:consent}
\end{enumerate}

\subsection*{Should Do}

Not launch-blocking.  Fast follow-up candidates.

\begin{enumerate}[nosep,leftmargin=2.5em]
  \featureitem{In-display agent}{Anthropic SDK, LuxExchange tool-use
    loop, leaf/composite commands, five-mode permission
    policy}\label{feat:agent}
  \featureitem{LLM extension authoring endpoint}{\texttt{/api/extensions/author}
    driven by in-display agent; screenshot-iterate
    loop}\label{feat:authoring}
  \featureitem{Community registry}{GitHub-hosted extension store with
    publish and search; ratings, download counts, verification badges
    (lint, CI, tests)}\label{feat:community}
  \featureitem{Reference extensions}{matplotlib plot, terminal
    (pyte), plotly snapshot, sankey diagram, confusion
    matrix}\label{feat:reference}
  \featureitem{Texture rendering mode}{extensions produce pixel
    buffers displayed as ImGui images; integrates matplotlib and
    PIL}\label{feat:texture}
  \featureitem{Window rendering mode}{extensions own native OS
    windows; Lux tracks lifecycle and coordinates
    focus}\label{feat:window}
  \featureitem{Remote transport}{TCP with TLS; multi-machine agent
    sessions connecting to one display}\label{feat:remote}
  \featureitem{Screenshot API}{capture display state as image for
    LLM-driven GUI iteration}\label{feat:screenshot}
\end{enumerate}

\subsection*{Won't Do}

Explicitly excluded to clarify scope.

\begin{enumerate}[nosep,leftmargin=2.5em]
  \featureitem{Web-based rendering}{Lux is Python-native.  HTML,
    Electron, and webviews are out of scope.  If content needs web
    embedding, use a web-native tool}\label{feat:no-web}
  \featureitem{GUI framework}{Lux hosts extensions; it does not
    define a widget class hierarchy.  No Lux-specific widget
    SDK}\label{feat:no-framework}
  \featureitem{Local LLM inference}{the in-display agent calls
    Anthropic's API; Lux does not host models
    locally}\label{feat:no-inference}
  \featureitem{Window manager}{OS-native windows from window-mode
    extensions are coordinated, not composited.  No tiling WM, no
    compositor}\label{feat:no-wm}
  \featureitem{Mobile or tablet}{desktop only (macOS,
    Linux)}\label{feat:no-mobile}
\end{enumerate}

% ══════════════════════════════════════════════════════════════════════════════
% REFERENCES
% ══════════════════════════════════════════════════════════════════════════════

\newpage
\printbibliography[title={References}]

\end{document}
